% !TeX root = ../tfg.tex
% !TeX encoding = utf8

% Recordatorio de inferencia general
% https://alejandroschuler.github.io/mci/introduction-to-modern-causal-inference.html


% TODO: perhaps move this to an appendix
\chapter{Preliminares}

\begin{definicion}[Población]
    Una población de estudio es un conjunto de ejemplares a estudiar.
\end{definicion}

\begin{definicion}[Espacio de probabilidad]
    $(\Omega, \mathcal{F}, P)$ es un espacio de probabilidad si
    \begin{enumerate}
        \item $\Omega$ es un conjunto no vacío.
        \item $\mathcal{F} \subseteq 2^\Omega = \{U | U \subseteq \Omega \}$ es una $\sigma$-álgebra.
        \item $P: \mathcal{F} \mapsto [0, 1]$ es una medida de probabilidad (medida tal que $P(\Omega) = 1$).
    \end{enumerate}
\end{definicion}

\begin{definicion}[Variable aleatoria]
    Se dice que $X:\Omega \mapsto E$ es una variable aleatoria si:
    \begin{enumerate}
        \item $(\Omega, \mathcal{F}, P)$ es un espacio de probabilidad.
        \item $E$ es un espacio medible (típicamente $\R^n$).
        \item $X$ es una función medible.
    \end{enumerate}
\end{definicion}

% TODO: source
\begin{definicion}[Muestra Aleatoria Simple, con reemplazamiento]\label{def:muestra}
    Una muestra de una variable aleatoria $X$ es una $n$-tupla $(X_1,...,X_n)$
    donde $X_i$ tiene la distribución de $X$ para cada $i \in \{1, ..., n\}$, y son
    mutuamente independientes.

    \begin{displaymath}
        \begin{cases}
            P(X_1 \in B_1, ..., X_n \in B_n) = \prod_{i=1}^n P(X_i \in B_i) \;\forall B_1, ..., B_n \text{ medibles} \\
            X_i \sim X \;\forall i = 1, ..., n
        \end{cases}
    \end{displaymath}

    Esto equivale a un vector de $n$ elementos de $P$ elegidos al azar
    con distribución uniforme sucesivamente y con reemplazamiento,
    justificando el nombre.
\end{definicion}

\begin{corolario}
    % TODO: fix space after MAS
    Dada una \MAS con reemplazamiento $(X_1, ..., X_n)$, se da que:
    \[
        P(X_1, ..., X_n) = \prod_{i=1}^n P(X_i)
    \]
\end{corolario}

% TODO: source
% TODO: maybe delete
\begin{definicion}[Muestra Aleatoria simple sin reemplazamiento]
    Una muestra $M$ donde los elementos de la población se eligen de forma que
    \[ \forall k \leq n, P(S_1) = P(S_2) \forall S_1, S_2 \in \{U \in 2^P | \#(U) = k \},  \]

    Donde $\#$ denota la función que cuenta el número de elementos de un conjunto.
    Esto es, la probabilidad de elegir cualquier subconjunto subconjunto con $k$ elementos es la misma
    que la de cualquier otro subconjunto de $k$ elementos.
\end{definicion}

% TODO: maybe delete, see above
\begin{corolario}
    Una muestra aleatoria simple sin reemplazo no es independiente, puesto que
    no se puede elegir el mismo sujeto de la población dos veces.
    % Sin embargo, sí conserva otras propiedades, como la intercambiabilidad, de la que hablaremos más tarde.
    % TODO: will we talk about it?
    % https://en.wikipedia.org/wiki/Exchangeable_random_variables
\end{corolario}

\begin{ejemplo}
    \begin{enumerate}
        \item Altura de una persona elegida al azar, medida en cm: $\{175\}$.
        \item Altura y peso de dos personas elegidas al azar con reemplazo, medida en cm y kilogramos:
              $\{(175, 75), (168, 66)\}$.
    \end{enumerate}
\end{ejemplo}

\begin{definicion}[Modelo]\cite[1.2]{mci}\cite[1.3.1]{targetedLearning}\label{def:modelo}
    Un modelo $\mathcal{M}$ es un par $(X, \mathcal{P})$ donde $X$ es una variable aleatoria
    y $\mathcal{P}$ es una familia de distribuciones de probabilidad.
\end{definicion}

\begin{nota} \label{nota:correctamenteEspecificado}
    Se dice que un modelo está \say{correctamente especificado} cuando incluye la verdadera distribución
    de probabilidad de la variable de estudio.
\end{nota}

\begin{notac}
    Denotaremos $P \in \mathcal{M}$ para indicar $P \in \mathcal{P}$, obviando el espacio muestral.
\end{notac}

\begin{observacion}
    Un modelo representa una colección de asunciones respecto a la generación de muestras.
\end{observacion}


% TODO: clear this up
\begin{definicion}[Estimando]
    Un estimando es una cantidad asociada a la distribución deseada de ser estimada.
    Formalmente, puede describirse como una función $\psi: \mathcal{M} \to \R^N$, donde $\mathcal{M}$ es un
    modelo, definido en \ref{def:modelo}.
\end{definicion}

\begin{nota}
    Un estimando puede tener como dominio tanto un modelo estadístico como un modelo causal. Nos referiremos a ellos
    respectivamente como estimandos estadísticos y causales.
\end{nota}

\begin{definicion}[Estadístico]
    % https://alejandroschuler.github.io/mci/52bed1d6fcf441baacee6b11048130f8.html
    Un estadístico lo definiremos como cualquier función (medible) de una muestra (\ref{def:muestra}) $\phi: \R^n \to \R^N$.
    % TODO: should I say the following? Con algo más de floritura y detalle, en cualquier caso.
\end{definicion}

\begin{definicion}[Estimador]
    Un estadístico $\hat\psi$ que pretende aproximar un estimando $\psi$.
    Si el estimando sólo toma valores en un subconjunto $\operatorname{Im}(\psi) = \Omega \subset \R^N$, típicamente pediremos que el estadístico
    tome valores en el mismo.
\end{definicion}

\begin{observacion}
    En la asignatura de Inferencia Estadística típicamente tomamos modelos parametrizados y comúnmente usamos
    estimadores de esos parámetros de una distribución, y por ello se exige
    que su imagen esté contenida en el espacio paramétrico.
\end{observacion}

En las próximas definiciones consideramos una variable aleatoria $X$ con distribución
$P$, una sucesión $\{X_i\}_{i \in \N}$ tal que $(X_1, \ldots, X_n)$ es una muestra
aleatoria simple de $X$ para todo $n \in \N$, y un estimador
$\hat\psi_n = \hat\psi(X_1, \ldots, X_n)$ de un estimando $\psi$.

\begin{definicion}[Estimador consistente]\label{def:estimador-consistente}
    En las condiciones anteriores, se dice que $\hat\psi_n$ es consistente
    si converge en probabilidad a $\psi(P)$. Eso es, lo es si para todo $\varepsilon > 0$
    \[
        \lim_{n\to\infty} P(|\hat\psi_n - \psi(P)| > \varepsilon) = 0
    \]
\end{definicion}

\begin{definicion}[Estimador asintóticamente insesgado]\label{def:insesgadez-asintotica}
    En las condiciones anteriores, se dice que $\hat\psi_n$ es asintóticamente insesgado si
    \[
        \lim_{n\to\infty}\E[\hat\psi_n] = \psi(P)
    \]
\end{definicion}

\begin{teorema}[La probabilidad total]\label{thm:probabilidad-total}
    Sea $(\Omega, \mathcal{F}, P)$ un espacio de probabilidad, $X$ una variable aleatoria
    con función de densidad $f_X$ y $A$ un evento en el espacio de probabilidad.
    Entonces,
    \[
        P(A) = \int_\R P(A | X = x)f_X(x)dx
    \]

    Si $X$ es discreta y $\{P_i\}_{i\in\N}$ es una partición de $\Omega$, entonces
    \[
        P(A) = \sum_i P(A | P_i) P(P_i)
    \]
\end{teorema}

\begin{definicion}[Función de pérdida]\cite[2.4]{elements_of_statistical_learning}
    Una función de pérdida $L: \R^2 \to \R$ mide la diferencia entre el valor predicho y el valor real.
\end{definicion}

\begin{ejemplo}[Función de pérdida cuadrática]
    La función de pérdida más común es la cuadrática:
    \[
        L(y, \hat{y}) = (y - \hat{y})^2
    \]
\end{ejemplo}

\begin{proposicion}[Predictor óptimo bajo pérdida cuadrática]
    Dado un vector aleatorio $(X, Y)$ con distribución conjunta $P_{X, Y}$, la función que
    minimiza el funcional $Err(f) := \E[(Y - f(X))^2]$ es
    \[
        f(x) = \E[Y | X = x]
    \]
    donde $E[Y | X = x]$ es la esperanza condicional de $Y$ dado $X = x$.
\end{proposicion}

\begin{definicion}[Problema de regresión] %\cite{targetedLearning}
    Dado un vector aleatorio $(X,Y)$, con distribución conjunta en $\mathcal{M}$,
    y con $X$ tomando valores en $\R^{d_x}$, e $Y$ tomando valores en $\R^{d_y}$,
    consideramos el siguiente estimando:
    \begin{align*}
        \Psi:   & \quad \mathcal{M} \to (\R^{d_x} \to \R^{d_y}) \\
        \Psi(P) & = \E_P[Y | X] : \R^{d_x} \to \R^{d_y}
    \end{align*}
    Que toma como valor:
    \[
        \Psi(P)(x) = \E_P[Y | X = x]
    \]
    El problema de regresión consiste en estimar esta función.
\end{definicion}

\begin{definicion}[Modelo de regresión]
    Dadas variables aleatorias $X$ e $Y$, un modelo de regresión consiste en una estimación de $Y$
    dado $X$. Se suele escribir funcionalmente como
    \[
        Y = f_\theta(X, A) + \epsilon
    \]

    Supongamos que nuestra variable de estudio $Y$ sigue una distribución desconocida pero contenida en
    un modelo $\mathcal{M}$ paramétrico, de forma que sigue una forma funcional $Y = f(X, A)$

    modelo paramétrico $\mathcal{M}_\theta, \theta \in \Theta$, de forma que
    $Y = f_\theta(X, A) + \epsilon$,
    para cierto $\theta \in \Theta$ y $\epsilon$ una variable aleatoria con distribución conocida.


    Comúnmente se asume $\E[\epsilon | X, A] = 0$ y homocedasticidad, es decir, $Var[\epsilon | X, A] = \sigma^2$ (ver \autoref{ap:hipotesis-regresion}).
\end{definicion}
